%%%%%
%
\section{Calcul simplifi� des 2 composantes du vent horizontal avec
reconnaissance de l'unit�}

\itemstyle{Op\'erateur~:} {\sf CalcVent\_h\_simple\_auto\_unit}
\index{CalcVent\_h\_simple\_auto\_unit}

%%%%%OBJET%%%%
\itemstyle{Objet~:}
Cet op�rateur calcule les 2 composantes du vecteur vent horizontal. 
Le calcul s'effectue avec les formules simplifi�es. Ces formules sont 
valides dans le cas o� l'angle d'incidence de l'air par rapport
au mobile, l'angle de tangage, et l'angle de roulis sont petits.\newline
Cet op�rateur travaille avec les diff�rents types d'unit�s dans la
mesure o� celles-ci restent homog�nes.Il faut lui passer en argument l'unit� de
mesure des angles.    


%%%%%ARGUMENTS%%%%
\bigskip
%%%%%ENTREES%%%%
\begin{tabarg}
\multicolumn{3}{l}{\titretab {Entr\'ees}}\\
\hline
$V_p$ & Vecteur, {\tt float} & Vitesse propre ($m/s$ ou $kt$)\\
$\beta$ & Vecteur, {\tt float} & Angle de d\'erapage ($radians$ ou $degres$)\\
$\psi$ & Vecteur, {\tt float} & Cap de la plate-forme mobile 
(en $radians$ ou $degres$, dans le sens indirect � partir du Nord g�ographique)\\
$unit$ & Coeff., {\tt String         } & Unite de mesure des angles ($radians$
ou $degres$)\\
$V_e$ & Vecteur, {\tt float} & Composante zonale de la vitesse de la plate-forme mobile 
par rapport au sol ($m/s$)\\
$V_n$ & Vecteur, {\tt float} & Composante m\'eridienne de la vitesse de la plate-forme mobile
par rapport au sol ($m/s$)\\
\end{tabarg}

\medskip
%%%%%SORTIES%%%%
\begin{tabarg}
\multicolumn{3}{l}{\titretab {Sorties}}\\
\hline
$U$ & Vecteur, {\tt float} & Composante zonale du vent ($m/s$)\\
$V$ & Vecteur, {\tt float} & Composante m\'eridienne du vent  ($m/s$)\\
\end{tabarg}


%%%%%FORMULE%%%%
\itemstyle{Formule}

%%CalcVent_h_simple_auto_unit
\begin{eqnarray*}
U & = & V_e - V_p . \sin (\psi + \beta) \\
V & = & V_n - V_p . \cos (\psi + \beta) 
\end{eqnarray*}



%%%%%AUTEUR%%%%
\itemstyle{Auteur~:}
CNRM/GMEI/TRAMM

%%%%%REFERENCE%%%%
\itemstyle{R\'ef\'erences~:}
Bulletin NCAR-RAF Nr 23 \cite{RAF-23}.
